% Figure: an open-bottom diving bell lowered to depth d. The trapped air
% compresses as the ambient water pressure rises, so the water rises inside
% the bell and the breathable air space shrinks. No em-dashes.
\begin{figure}[htbp]
\centering
\begin{tikzpicture}[scale=1.0]
% free surface
\draw[engblue, thick] (-0.5,6.0) -- (9.5,6.0);
\node[engblue, font=\scriptsize, anchor=south] at (-0.3,6.0) {free surface};
% water body shading
\fill[engblue!10] (-0.5,0) rectangle (9.5,6.0);
% Bell at the surface (left), full of air
% bell wall (open bottom): inverted U
\draw[engblack, very thick] (0.6,5.4) -- (0.6,4.0) -- (2.4,4.0) -- (2.4,5.4);
\draw[engblack, very thick] (0.6,5.4) -- (2.4,5.4);
% air fills the bell to the rim
\fill[engamber!25] (0.62,4.02) rectangle (2.38,5.38);
\node[engamber!80!black, font=\scriptsize] at (1.5,4.7) {air};
\node[engblack, font=\scriptsize, anchor=north] at (1.5,4.0) {at surface};
% Bell at depth (right): water risen inside
\draw[engblack, very thick] (6.6,3.0) -- (6.6,1.6) -- (8.4,1.6) -- (8.4,3.0);
\draw[engblack, very thick] (6.6,3.0) -- (8.4,3.0);
% compressed air near the top, water below
\fill[engamber!25] (6.62,2.55) rectangle (8.38,2.98);
\node[engamber!80!black, font=\scriptsize] at (7.5,2.76) {air};
\fill[engblue!30] (6.62,1.62) rectangle (8.38,2.55);
\node[engblue!70!black, font=\scriptsize] at (7.5,2.05) {risen water};
\node[engblack, font=\scriptsize, anchor=north] at (7.5,1.6) {at depth $d$};
% depth arrow to the submerged bell rim
\draw[enggray, dashed] (8.4,6.0) -- (9.2,6.0);
\draw[enggray, dashed] (8.4,2.98) -- (9.2,2.98);
\draw[enggray, <->] (9.2,6.0) -- (9.2,2.98);
\node[enggray, font=\scriptsize, anchor=west] at (9.2,4.5) {$d$};
% descent arrow between the two bells
\draw[engred, ultra thick, ->] (2.9,4.4) .. controls (4.6,3.6) .. (6.2,2.9);
\node[engred, font=\scriptsize] at (4.6,3.9) {descent};
\end{tikzpicture}
\caption[Diving bell air pocket under compression]{An open-bottom diving
bell traps a fixed mass of air. At the surface the air fills the bell to
the rim. Lowered to depth $d$, the trapped air feels the raised ambient
pressure $p_{\mathrm{atm}}+\rho g d$ at the rim and compresses by Boyle's
law, so water rises inside the bell and the breathable space shrinks. The
air volume follows $p_1V_1=p_2V_2$ with the pressures taken at the trapped
gas, and the interface settles where the internal air pressure matches the
water pressure at the internal air-water surface. Pumping fresh air down
to restore the space is what a working bell or a pneumatic caisson does to
keep the chamber open at depth.}
\label{fig:vol03ch10:diving-bell}
\end{figure}
